%% Template feito com base no AbnTex2 para o Programa de Pós Graduação em Antropologia da UFMG
%%
%% This work may be distributed and/or modified under the conditions of the LaTeX Project Public License, either version 1.3 of this license or (at your option) any later version. The latest version of this license is in http://www.latex-project.org/lppl.txt and version 1.3 or later is part of all distributions of LaTeX version 2005/12/01 or later. This work has the LPPL maintenance status `maintained'. Further information about abnTeX2 are available on https://github.com/abntex/abntex2 | https://www.abntex.net.br/
%%
%% The Current maintainer of this template is Geraldo Pereira de M. Jr. 
%% Contact: geraldo.pmj@gmail.com
%%
\documentclass[
	12pt,				% tamanho da fonte
	openright,			% capítulos começam em pág ímpar (insere página vazia caso preciso)
	oneside,			% para impressão em recto e verso. Oposto a oneside
	a4paper,			% tamanho do papel.
	sumario = tradicional,
	% -- opções da classe abntex2 --
	%chapter=TITLE,		% títulos de capítulos convertidos em letras maiúsculas
	%section=TITLE,		% títulos de seções convertidos em letras maiúsculas
	%subsection=TITLE,	% títulos de subseções convertidos em letras maiúsculas
	%subsubsection=TITLE,% títulos de subsubseções convertidos em letras maiúsculas
	% -- opções do pacote babel --
	english,			% idioma adicional para hifenização
	french,				% idioma adicional para hifenização
	spanish,			% idioma adicional para hifenização
	brazil				% o último idioma é o principal do documento
	]{abntex2}
\input{Documento/config/packages}
\input{Documento/config/pdfconfig}

% ---
% Pacote que deixa vários comandos do LaTex comumente usados em português. Pode ver exemplos na Introdução. Contudo pode continuar a usar os comandos do LaTex normalmente se preferir.
% ---
\usepackage{PTLatexCommands}

% ---
% Informações de dados para CAPA e FOLHA DE ROSTO
% ---
\titulo{Título do trabalho}
\autor{Autor(a)}
\local{Belo Horizonte}
\data{\anoatual}
\orientador{Orientador(a)}
%\coorientador{Equipe \abnTeX}
\instituicao{%
  Universidade Federal de Minas Gerais
  \par
  Faculdade de Filosofia e ciências Humanas
  \par
  Programa de Pós-Graduação em Antropologia}
\tipotrabalho{Dissertação (Mestrado)}
% O preambulo deve conter o tipo do trabalho, o objetivo, 
% o nome da instituição e a área de concentração 
\preambulo{Dissertação de Mestrado apresentada ao
Programa de Pós-Graduação em Antropologia
da Faculdade de Filosofia e Ciências Humanas,
Universidade Federal de Minas Gerais como
requisito parcial ao título de mestre .\linebreak
\textbf{Área de concentração}: Arqueologia \newline
\textbf{Linha de Pesquisa}: Arqueologia do Mundo Moderno \newline
\textbf{Orientador}: \imprimirorientador}

% ---
% compila o indice
% ---
\makeindex
% ---

% ----
% Início do documento
% ----
\begin{document}
\selectlanguage{brazil}
% Retira espaço extra obsoleto entre as frases.
\frenchspacing 
%spaçamento 
%----------------------------------------------------
% ELEMENTOS PRÉ-TEXTUAIS
%----------------------------------------------------
\pretextual
%---
% Capa
%---
\imprimircapa 
%---
% Folha de rosto
% (o * indica que haverá a ficha bibliográfica)
%---

\imprimirfolhaderosto*

%---
% Inserir a ficha bibliografica
%---
% Isto é um exemplo de Ficha Catalográfica, ou ``Dados internacionais de
% catalogação-na-publicação''. Você pode utilizar este modelo como referência. 
% Porém, provavelmente a biblioteca da sua universidade lhe fornecerá um PDF
% com a ficha catalográfica definitiva após a defesa do trabalho. Quando estiver
% com o documento, salve-o como PDF no diretório do seu projeto e substitua todo
% o conteúdo de implementação deste arquivo pelo comando abaixo:
\input{capitulos/ficha.tex}
%---
%Inserir errata
%---
%\input{capitulos/errata.tex}
%---
% Inserir folha de aprovação
%---
\input{capitulos/aprov.tex} 
%---
% Dedicatória
%---
\input{capitulos/dedi.tex}
%---
% RESUMOS
%---
\input{capitulos/resumos.tex}
%---
%LISTAS
%---
\input{capitulos/Listas.tex}
%---
%inserir o sumario
%---
\input{capitulos/sum.tex}

%-----------------------------------------------------
% ELEMENTOS TEXTUAIS
%-----------------------------------------------------

\textual

%-----------------------------------------------------
%-----------------------------------------------------
\input{capitulos/introducao.tex}
%-----------------------------------------------------
\input{capitulos/cap1.tex}
%-----------------------------------------------------
\input{capitulos/cap2.tex}
%-----------------------------------------------------
\input{capitulos/cap3.tex}
%-----------------------------------------------------
\input{capitulos/cap4.tex}
%-----------------------------------------------------
\input{capitulos/conclusao.tex}
%-----------------------------------------------------

% ELEMENTOS PÓS-TEXTUAIS ------------------------------------------------------

\postextual

\pagestyle{plain}

%----------------------------------------------------
% Referências bibliográficas
%----------------------------------------------------

\printbibliography[title={Referências Bibliográficas}]

%-----------------------------------------------------
% Glossário 
%-----------------------------------------------------
% Consulte o manual da classe abntex2 para orientações sobre o glossário.

%\glossary

%-----------------------------------------------------
% Apêndices 
%-----------------------------------------------------
% Inicia os apêndices
%---

\begin{apendicesenv}
\input{capitulos/apend.tex}
\end{apendicesenv}

%--- 

%-----------------------------------------------------
% Anexos
%-----------------------------------------------------

\begin{anexosenv}
\input{capitulos/anex.tex}
\end{anexosenv}

%-----------------------------------------------------
% INDICE REMISSIVO
%-----------------------------------------------------

\input{capitulos/indiceremiss.tex}

%-----------------------------------------------------
\end{document}
